% Generated draft: requires mathematical and source review.
\documentclass[11pt]{article}
\usepackage[margin=25mm]{geometry}
\usepackage{fontspec}
\setmainfont{Latin Modern Roman}
\usepackage{amsmath,amssymb,mathtools,unicode-math}
\setmathfont{Latin Modern Math}
\usepackage{enumitem,needspace,xurl,hyperref}
\setlength{\parindent}{0pt}
\setlength{\parskip}{5pt}
\setlength{\emergencystretch}{3em}
\newcommand{\N}{\mathbb{N}}
\newcommand{\Z}{\mathbb{Z}}
\newcommand{\Q}{\mathbb{Q}}
\newcommand{\R}{\mathbb{R}}
\newcommand{\C}{\mathbb{C}}
\newcommand{\F}{\mathbb{F}}
\newcommand{\A}{\mathbb{A}}
\newcommand{\PP}{\mathbb{P}}
\newcommand{\E}{\mathbb{E}}
\newcommand{\eps}{\varepsilon}
\newcommand{\dd}{\,\mathrm{d}}
\begin{document}
\section*{UnsolvedMath ID: 10600005}
\textbf{Source problem code:} catalogue entry
\subsection{Definitions and mathematical statement}
The compound item concerns a proposed theory of virtual three-manifolds, defined as equivalence classes of virtual framed-link diagrams under generalized Kirby moves. Let \(\mathcal D_{\mathrm{cl}}\) and \(\mathcal D_{\mathrm{vir}}\) denote the classical and virtual diagram sets, and let \(\sim_K\) and \(\sim_{vK}\) denote the corresponding Kirby relations. The precise generalized move list is not supplied.
\begin{enumerate}
\item Do homeomorphism and virtual equivalence coincide on ordinary three-manifolds? Explicitly, for every pair of classical framed surgery links \(K,L\), does
\[K\sim_{vK}L\quad\Longrightarrow\quad K\sim_K L?\]
The usual classical surgery/Kirby conventions identify the latter with homeomorphism of the presented manifolds in the intended category. Equivalently, assuming the virtual move system includes classical Kirby moves, is the induced map
\[\mathcal D_{\mathrm{cl}}/\!\sim_K\ \longrightarrow\ \mathcal D_{\mathrm{vir}}/\!\sim_{vK}\]
injective? Intermediate diagrams in \(\sim_{vK}\) may be virtual even though the endpoints are classical.
\item Give a geometric-topological interpretation of the virtual equivalence classes. A solution must identify geometric objects and an equivalence relation realizing these classes, rather than merely restating their formal diagram definition; the excerpt prescribes no particular geometric category.
\item Construct another virtual-three-manifold theory by surgery on links in thickened surfaces \(S_g\times\R\), considered up to stabilization, and determine whether it coincides with the theory attributed to Kauffman and Dye. One must specify the allowed links and framings, surface stabilizations, surgery and equivalence of the resulting objects, and a comparison of the two resulting class sets compatible with the intended presentations.
\end{enumerate}
Here \(S_g\) denotes the surface of genus \(g\) in the source's conventions. Its orientation/closedness, the treatment of the noncompact ends of \(S_g\times\R\), and how stabilization acts before and after surgery are unstated. No equivalence between the two theories, geometric realization or absence of new identifications of classical manifolds is assumed.
\subsection{Short English statement}
The item asks three things: whether virtual Kirby moves identify any additional classical surgery manifolds; what virtual three-manifolds mean geometrically; and whether surgery in stabilized thickened surfaces gives the Kauffman--Dye theory.
\subsection{Sources}
\begin{itemize}
\item Local UnsolvedMath registry, permanent ID \(10600005\), exact source code \texttt{catalogue entry}.
\item \emph{Unsolved Problems in Virtual Knot Theory and Combinatorial Knot Theory}, revised survey, problem-list item 5: \url{https://arxiv.org/abs/1409.2823}. Registry access 2026-07-29; source-tex extraction. The source attributes the comparison theory to Kauffman and Dye but its embedded reference is missing.
\item The research summary mentions an unverified lead titled \emph{Complexity of virtual 3-manifolds}, tentatively ``Martelli?'' and ``arXiv:1609.xxxx'' (2016). This is incomplete bibliographic data, not a verified citation or identifier.
\end{itemize}
\end{document}
